\section{Aritmética cardinal básica}

É possível definir operações de aritmética cardinal para simplificar o cálculo de cardinalidades infinitas.
Tais operações, por ora de soma e produto, não coincidem com a noção de aritmética ordinal.

\begin{definition}
    Sejam $\kappa$ e $\lambda$ cardinais.
    Define-se a soma cardinal $\kappa + \lambda$ como $|\{0\}\times \kappa \cup \{1\}\times \lambda|$.
    Define-se o produto cardinal $\kappa \cdot \lambda$ como $|\kappa \times \lambda|$.
\end{definition}

A definição acima faz sentido, uma vez que os conjuntos $\{0\}\times \kappa \cup \{1\}\times \lambda$ e $\kappa \times \lambda$ são bem ordenáveis, por exemplo, pela ordem lexicográfica ou pela ordem $\triangleleft$ da Seção~\ref{sec:quadradosDeOrdinais}.

\begin{example}
    Considere $\kappa=\lambda=\omega$.
    
    A \textbf{soma ordinal} $\omega+\omega$ é o ordinal $\omega + \omega>\omega$, um ordinal enumerável.
    O \textbf{produto ordinal} de $\omega\cdot\omega$ é o ordinal $\omega \cdot \omega>\omega+\omega>\omega$, um ordinal enumerável.

    A \textbf{soma cardinal} de $\omega+\omega$ é $|\{0\}\times \omega \cup \{1\}\times \omega|=|\omega|=\omega$.
    O \textbf{produto cardinal} de $\omega$ e $\omega$ é $|\omega \times \omega|=|\omega|=\omega$.
\end{example}

Na literatura, as duas notações são utilizadas sem aviso.
A aritmética cardinal é de uso muito mais comum do que a aritmética ordinal, de modo que, quase sempre, as operações de soma e produto são operações de aritmética cardinal.
Como o uso de aritmética ordinal é mais restrito, muitas vezes fica claro pelo contexto quando as operações de soma e produto são operações de aritmética ordinal.
Quando não fica claro, tal ambiguidade é usualmente resolvida por meio de um esclarecimento por escrito.
Em última instância, algum símbolo diferente para um dos pares de operações é pontualmente adotado, mas isso é raro de ser encontrado na literatura.

Didaticamente, ao utilizarmos letras $\kappa$, $\lambda$ e $\mu$ e $|A|$, bem como alephs, estamos, a menos que dito o contrário, nos referindo a cardinais, e, portanto, as operações de soma e produto que utilizaremos serão operações de aritmética cardinal.

Já ao utilizarmos letras como $\alpha$, $\beta$ e $\gamma$ e $\delta$, estamos, a menos que dito o contrário, nos referindo a ordinais, e, portanto, as operações de soma e produto que utilizaremos serão operações de aritmética ordinal.

Antes de explanarmos como operar aritmética cardinal de forma prática, estudaremos lemas que a tornam extremamente útil.

\begin{lemma}
    Seja $A$ um conjunto e suponha que exista um conjunto bem ordenável $B$ e $f:B\to A$ uma função sobrejetora.
    Então $A$ é bem ordenável e $|A|\leq |B|$.
\end{lemma}
\begin{proof}
    Seja $<_B$ uma boa ordem em $B$.

    Defina $g: A\to B$ dada por $g(a)=\min\{ b \in B: f(b)=a\}$.

    $g$ é claramente injetora.
    Defina uma relação $<_A$ em $A$ por:
    \[
    x <_A y \iff g(x) <_B g(y).
    \]
    Note que $g: (A, <_A) \to (g[A], <_B)$ é um isomorfismo de ordens, e, portanto, $A$ é bem ordenável.
\end{proof}

\begin{lemma}
Sejam $A$ e $B$ conjuntos bem ordenáveis.
Então $A\times B$ é bem ordenável e $|A\times B|=|A|\cdot |B|$.
\end{lemma}
\begin{proof}
    Sejam $\kappa=|A|$ e $\lambda=|B|$.
    Sejam $f_A: \kappa \to A$ e $f_B: \lambda \to B$ funções bijetoras.
    Seja $g: \kappa \times \lambda \to A\times B$ dada por $g(\alpha, \beta) = (f_A(\alpha), f_B(\beta))$.
    É imediato que $g$ é uma função bijetora, e, portanto, $A\times B$ é bem ordenável e $|A\times B|=|\kappa \times \lambda|=\kappa \cdot \lambda = |A|\cdot |B|$.
\end{proof}

\begin{lemma}
Sejam $A$ e $B$ conjuntos bem ordenáveis disjuntos.
Então $A\cup B$ é bem ordenável e $|A\cup B|=|A|+|B|$.
\end{lemma}
\begin{proof}
    Suponha que $A$ e $B$ são disjuntos.
    Temos que $A\approx(\{0\}\times A)$ e $B\approx(\{1\}\times B)$.
    Assim, $A\cup B \approx (\{0\}\times A)\cup (\{1\}\times B)$.
    Portanto, $A\cup B$ é bem ordenável e $|A\cup B|=|A|+|B|$.
\end{proof}

Operacionalmente, temos a seguintes preservações de desigualdades.
\begin{lemma}
Sejam $\kappa$, $\kappa'$, $\lambda$ e $\lambda'$ cardinais tais que $\kappa \leq \kappa'$ e $\lambda \leq \lambda'$.
Então $\kappa + \lambda \leq \kappa' + \lambda'$ e $\kappa \cdot \lambda \leq \kappa' \cdot \lambda'$.
\end{lemma}
\begin{proof}
    Para a primeira desigualdade, note que $\{0\}\times \kappa \cup \{1\}\times \lambda \subseteq \{0\}\times \kappa' \cup \{1\}\times \lambda'$.
    Assim,
    \[
    \kappa + \lambda = |\{0\}\times \kappa \cup \{1\}\times \lambda| \leq |\{0\}\times \kappa' \cup \{1\}\times \lambda'| = \kappa' + \lambda'.
    \]
    Para a segunda desigualdade, note que $\kappa \times \lambda \subseteq \kappa' \times \lambda'$.
    Assim,
    \[
    \kappa \cdot \lambda = |\kappa \times \lambda| \leq |\kappa' \times \lambda'| = \kappa' \cdot \lambda'.
    \]
\end{proof}

Sem assumir a disjunção de $A$ e $B$, temos a seguinte desigualdade.
\begin{lemma}
Sejam $A$ e $B$ conjuntos bem ordenáveis.
Então $A\cup B$ é bem ordenável e $|A\cup B|\leq |A|+|B|$.
\end{lemma}
\begin{proof}
    Seja $B'=B\setminus A$.
    Temos que $A$ e $B'$ são bem ordenáveis e disjuntos, e, portanto, $A\cup B'=A\cup B$ é bem ordenável e disjunto.
    Além disso, $|B'|\leq |B|$.
    Portanto:
    \[
    |A\cup B| = |A\cup B'| = |A| + |B'| \leq |A| + |B|.
    \]
\end{proof}

Agora mostraremos as principais ferramentas para cálculos envolvendo aritmética cardinal.

\begin{lemma}\label{lem:aritmeticaCardinalComutativa}
    Sejam $\kappa$ e $\lambda$ cardinais tais que $2\leq \kappa$ e $2\leq \lambda$.
    Então $\kappa + \lambda=\lambda+\kappa$ e $\kappa \cdot \lambda=\lambda \cdot \kappa$.
\end{lemma}
\begin{proof}
    A prova é simples e deixada como exercício ao leitor.
\end{proof}

\begin{lemma}\label{lem:aritmeticaCardinalBasica}
    Sejam $\kappa$ e $\lambda$ cardinais tais que $2\leq \kappa$ e $2\leq \lambda$.
    Então $\kappa + \lambda\leq \kappa \cdot \lambda$.
\end{lemma}
\begin{proof}
    Seja $\mu=\max\{\kappa, \lambda\}$.
    Se $\mu=\lambda$, então
    \[
    \{0\}\times \kappa \cup \{1\}\times \lambda \subseteq \{0\}\times \kappa \cup \{1\}\times \lambda\subseteq \kappa \times \lambda.
    \]

    Portanto, $\kappa+\lambda\leq \kappa\cdot \lambda$.

    Já se $\mu=\kappa$, pelo argumento anterior e pelo Lema~\ref{lem:aritmeticaCardinalComutativa}, temos que:

    \[
    \kappa+\lambda = \lambda + \kappa \leq \lambda \cdot \kappa = \kappa \cdot \lambda.
    \]
\end{proof}

\begin{theorem}
    Sejam $\kappa$ e $\lambda$ cardinais infinitos.
    Então $\kappa + \lambda = \kappa \cdot \lambda = \max\{\kappa, \lambda\}$.
\end{theorem}
\begin{proof}
    Seja $\mu=\max\{\kappa, \lambda\}$.
    É imediato que $\mu\leq \kappa + \lambda$, uma vez que $\mu$ injeta em $\{0\}\times \kappa \cup \{1\}\times \lambda$.

    Assim:

    \[
    \mu \leq \kappa + \lambda \leq \kappa \cdot \lambda \leq \mu \cdot \mu = \mu.
    \]

    A última desigualdade é verdadeira pelo Corolário~\ref{cor:quadradoDeCardinalInfinito}.

    Portanto, $\kappa + \lambda = \kappa \cdot \lambda = \mu=\max\{\kappa, \lambda\}$.
\end{proof}
Operar cardinais finitos com infinitos também é uma tarefa simples.
\begin{lemma}
    Seja $n$ um cardinal finito e $\kappa$ um cardinal infinito.
    Então:
    \begin{enumerate}[label=(\alph*)]
        \item $n+\kappa = \kappa + n = \kappa$;
        \item $n\cdot \kappa = \kappa \cdot n = \kappa$ se $n\neq 0$;
        \item $0\cdot \kappa = \kappa \cdot 0 = 0$.
    \end{enumerate}
\end{lemma}
\begin{proof}
    Para (a), note que $\kappa\leq n+\kappa\leq \kappa+\kappa = \kappa$.
    Assim, $n+\kappa = \kappa + n = \kappa$.

    Para (b), note que $\kappa\leq n\cdot \kappa\leq \kappa\cdot \kappa = \kappa$.
    Assim, $n\cdot \kappa = \kappa \cdot n = \kappa$.

    Para (c), note que $0\cdot \kappa =|0\times \kappa|=|\emptyset|=0$.
\end{proof}

Finalmente, a aritmética cardinal finita coincide com a usual (e, portanto, com a aritmética ordinal).

\begin{lemma}
    Sejam $m$ e $n$ cardinais finitos (ou seja, números naturais).
    Então os cardinais $m+n$ e $m\cdot n$ obtidos pelas definições de aritmética cardinal coincidem com os cardinais $m+n$ e $m\cdot n$ obtidos pelas definições usuais de soma e produto de números naturais.
\end{lemma}

\begin{proof}
    Sejam $m$ e $n$ cardinais finitos.
    Então $m\approx \{0\}\times m$ e $n\approx \{0\}\times n$.

    Pelo Exercício~\ref{ex:aritmeticaFinitaUniao} o número $m+n$ obtido pela aritmética cardinal usual coincide com $|\{0\}\times m \cup \{1\}\times n|$, e, portanto, com o número $m+n$ obtido pela aritmética cardinal.

    Analogamente, pelo Exercício~\ref{ex:aritmeticaFinitaProduto} o número $m\cdot n$ obtido pela aritmética cardinal usual coincide com $|m\times n|$, e, portanto, com o número $m\cdot n$ obtido pela aritmética cardinal.
\end{proof}

\subsection*{Exercícios}
\begin{exercise}Prove o Lema~\ref{lem:aritmeticaCardinalComutativa}.
\end{exercise}